\section{鸽巢原理与必然碰撞}
\label{sec:04-Discrete-Mathematics/01-Combinatorics/03}

你参加一个 13 人的聚会。有人打赌说，这群人里至少有两个人的生日在同一个月份。你不必挨个问每个人的生日就可以确信他赢——12 个月装 13 个人，总有一个月至少装了两个。

这个推理简单到你可能觉得它算不上一个原理。但问题不在这里。问题在于：当一个真实的数学题摆在你面前时，你是那个能一眼认出``鸽子''和``鸽巢''的人吗？

\subsection{数量超过容量，碰撞就不是偶然}

把 \(N\) 个对象放进 \(k\) 个盒子。如果每个盒子都尽量``客气''地少装——每盒至多装 \(\lceil N/k\rceil-1\) 个——那么总容量至多是

\[
k\left(\left\lceil\frac Nk\right\rceil-1\right)
<k\cdot\frac Nk=N.
\]

装不下。所以至少有一个盒子装了至少 \(\lceil N/k\rceil\) 个。这就是广义鸽巢原理的全部证明：容量不够迫使拥挤发生。

最简单的形式是 \(N>k\)：对象比盒子多，至少有一个盒子装了两个或更多。翻译成函数的语言：从大小为 \(N\) 的集合到大小为 \(k\) 的集合的函数，若 \(N>k\)，不可能是单射——总有两个不同输入被映射到同一个输出。

原理本身不需要更多解释了。真正值得花时间的，是练习在陌生的题目里认出这两样东西：谁是鸽子，谁是鸽巢，以及同盒意味着什么。

\subsection{一个余数例子}

随便取 6 个整数：\(7, 13, 22, 31, 45, 59\)。除以 5，余数分别是 \(2, 3, 2, 1, 0, 4\)。余数 2 出现了两次。不是巧合：6 个整数，余数只有 5 种可能（0 到 4），鸽子比巢多。

把这个推成一般形式：任取 \(n+1\) 个整数，至少有两个除以 \(n\) 的余数相同。鸽子是这些整数，鸽巢是 \(n\) 个余数类。同盒意味着什么？两个数的余数相同，它们的差就能被 \(n\) 整除。

你发现了吗，鸽巢原理在这个证明里做了一件很特别的事：它保证存在一对差能被 \(n\) 整除的整数，但不告诉你具体是哪一对。想找出这一对，你得真的去算每个数的余数、比较、配对——那是另一个问题了。原理只负责说``一定有''，不管``在哪里''。这是非构造性存在证明的标志：容量矛盾强迫某个结构必须出现，但你手里没有构造它的配方。

\subsection{相邻距离的例子}

在单位区间 \([0,1]\) 上任取 \(n+1\) 个点。把区间等分成 \(n\) 段，每段长 \(1/n\)。\(n+1\) 个点落入 \(n\) 个小区间，至少两个点落在同一段里——它们之间的距离不会超过 \(1/n\)。

具体一点：\(n=5\)，在 \([0,1]\) 上任取 6 个点。把区间分成 \([0,0.2), [0.2,0.4), [0.4,0.6), [0.6,0.8), [0.8,1]\)。6 个点，5 个区间，必有两个点在同一段里，距离不超过 0.2。

分点落到区间边界上的情形需要统一处理一下：如果刚好有一个点落在分界的 0.2 上，它该归左边的区间还是右边的？严谨的做法是规定前 \(n-1\) 段左闭右开，最后一段两端都闭合。这不改变思路，却影响证明的严密性——边界处理不当，一个点可能被同时塞进两个盒子，或者被漏掉。

\subsection{极值形式与反证}

广义鸽巢原理常常用反证法包装成更顺手的形状：假设每个盒子最多装 \(q\) 个，那么 \(k\) 个盒子最多装 \(kq\) 个。一旦实际数量超过 \(kq\)，反证假设就站不住了——至少有一个盒子装了至少 \(q+1\) 个。

用这个形式算一个实用的数字：100 个人，12 个月。若每月至多 8 人，总数至多 96。但实际有 100 人，所以至少有一个月有至少 9 人生日。结论是``至少 9''，不是``每个月都大约 8 个人''，也不是``一定有 9 个人都在一月''。原理只保证下界，不描述分布。

这个极值思考方式也揭示了鸽巢原理和``最坏情况分析''之间的关系：你想象一个对手，他拼命想让每个盒子尽量平均、尽量不拥挤。你算出他在这个策略下最多能装多少。如果你的对象比他的极限容量还多，那拥挤就不可避免了——连最坏情况都防不住。

\subsection{原理只保证存在，不保证均匀}

鸽巢原理给的是最坏情况下的下界。它不说明：

\begin{itemize}
\tightlist
\item
  哪些盒子拥挤；
\item
  分布是否接近均匀；
\item
  如何高效地找到那个拥挤的盒子。
\end{itemize}

这三件事分别需要不同的工具。如果你想找具体哪个盒子拥挤，需要算法。如果你想知道分布大概长什么样，可能需要概率——给对象一个随机分配模型，用期望或高概率结论来刻画。但鸽巢原理的优势恰好在于它不需要随机模型：它是纯确定性的，不依赖任何分布假设，因此在最苛刻的条件下仍然成立。

这个``纯确定性''也是它的代价：结论较弱。概率方法可能告诉你``大概率有一盒装了远比 \(\lceil N/k\rceil\) 多的对象''，而鸽巢原理只能说``至少有一盒装了 \(\lceil N/k\rceil\)''。两者的保证强度不同，适用的场景也不同。当概率模型可靠时，优先用概率工具；当没有任何随机假设可依靠时，鸽巢原理就是为数不多还能工作的武器。
